\documentclass[a4paper]{article}

\usepackage{graphicx}
\usepackage{amsmath,amssymb}
\usepackage{minted}
\usepackage{multirow}
\usepackage{float}
\usepackage{todonotes}
\usepackage{forest}
\usepackage{xstring}
\usepackage{xcolor}
\usepackage[most]{tcolorbox}
\usepackage{xparse}
\usepackage{natbib}

\usetikzlibrary{graphs,graphdrawing}
\usegdlibrary{layered}

% for external graphviz
\input{/home/donkarlo/Dropbox/repo/nd_ai_project/src/nd_ai/knoweldge_representation/language/formal/computer/programming/tex/latex/code_clips/external_graphviz.tex}

% for tags
\input{/home/donkarlo/Dropbox/repo/nd_ai_project/src/nd_ai/knoweldge_representation/language/formal/computer/programming/tex/latex/parts/control_sequence/kind/semtag/semtag.tex}

% for links
\input{/home/donkarlo/Dropbox/repo/nd_ai_project/src/nd_ai/knoweldge_representation/language/formal/computer/programming/tex/latex/code_clips/seehere.tex}

\title{}
\date{}
\author{Mohammad Rahmani}

\begin{document}
    \maketitle
    
    
    \section{Thema: Sollten Menschen weniger soziale Medien benutzen?}
        
        Schreibe etwa 80--100 Wörter. Du kannst diese Punkte benutzen:
        
        \begin{itemize}
            \item Wie oft benutzt du soziale Medien?
            \item Welche Vorteile haben sie?
            \item Welche Nachteile gibt es?
            \item Was ist deine persönliche Meinung?
            \item Was würdest du anderen Menschen empfehlen?
        \end{itemize}
        
        Persönlich benutze ich viel soziale Medien.
        
        Ich denke, dass die sozialen Medien sehr viel sinnvoll für die Enwicklung der Demokratie sein können.
        
        Ich weiß, dass wir bisher(so far) es nicht bezeugt(to withness) haben, aber wir müssen auch beachten, dass(dass spielst der Objektroll weil beaten brauchst ein davon(yeki az an=yek objekt)), obwohl diese Medien nicht sehr jung sind, sie auch nicht genug erwachsene(Adult) sind, damit wir erwarten, dass sie über unsere democratikche Zukunft ein effektiv Effekt zu haben.
        Aber, trotz allem, was ich gessagt habe, bezeugen(to witness) wir, dass diese Medien den stimmlosen Völker ein Mittle, um sich auszudrüken.
        
        Da ich denke, dass jedes Massenmedium(singular neutral), das(pronom für Massenmedien) der Gesellschaft(f) hinzufügt wird (von Behörde), dem Volk hift, durch Narrichten besser infrormiert zu werden.
        
        Nun hat sich die Frage in die verwandelt(to change, benutzt immer mit "in". Es heißt wir benutzen immer verwandeln in etwas die hier wir "in" dahin benutzt), wie die Narrichten uns helfen, auf eine bessere Demokratie(f) zuzumarschieren (Die zweite zu ist die zu die ist für infinitiv bauen benutzt worden).
        
        Nachrichten helfen uns, uns("uns" ist eine Variante von Refelxivpronomen für sich informieren die später in deselb Satz benutzt wird) besser über die Strecke(m, dativ Objekt) die, wir nehmen müssen, zu informieren (Das verb ist "sich informieren" aber "uns" ist vor ),  um unsere Rechte(plural hier aber Recht is neutral) besser zu verteidigen.
        
        
        Sogar vor(before, andres als bevor die Hauptsatz-nebensatz Konstruktion nötigt, vor muss immer vor ein Nom kommen. Weil es kein Bewegung gibt, dann Mann muss Dativ benuzen) der Strecke, die wir nehmen müssen, um unsere Rechte zu verteidigen, müssen wir über sie uns informieren.
        
        Außerdem müssen wir entdecken(to discover), welche politische Partei uns hilft (hilfen verb steht sich am ende, weil es "Welche" hat ein neben satz gebaut), (kein um...zu benutz wird, weil das "Helfen" konstruktion folgt "jemand Hefen zu etwas machen". Wir müssen wissen, obwohl was in die letze Phrase erklärt worden ist man kann helfen mit um ... zu benutzen aber die Bedutung wird ein bischen verschiden) unsere Rechte besser zu verteideigen.
        
        Wie es bisher erklärt worden ist(Nebenstaz), brauchen wir (zweichendrin die ) die richtigen Nachrichten(feminin plural), um eine praktische DemoKratie(f) zu haben.
        
        Wir können entweder ein einziges(wegen ein und Medium, wir brauchen ein gemischte Adjektiv, deshalb wir -es hinzufügen) Medium,(Neutrum Akkusativ Singular) von dem wir richtigesten Nachrichten(Akkusativobjekt von erwarten) erwarten, oder mehrere Medien von denen (Dativ plural), das Publikum die Nachrichten vergleichen kann, um die präzisesten Nachrichten(Akkusativ plural für entnehmen) zu entnehmen(to extract).
        
        Ich verteidige mehrere Medien, weil ich finde überhaupt nicht möglisch, dass wir von ein einzige sozial Medium erwarten, dass es ganz neutral bleibt.
        
        Deshlab(Deshalb macht nicht der Satz ein Nebensatz) finde ich mehrere Medien als die Quellen, von denen(Dativ weil es nachg von stelt und feminin, plural weil es ist ein Pronom für Quelen) wir über unsere Rechte(die Plural ist Rechte und der Nom Konjugirt nicht mit andren Teile des Satz ) und die politische Partei, die unsere Rechte verteidigt lernen können.
        
        
        Soziale Medien(pl) verbreiten(AKK obj) uns die außergewöhnlichen(Schwach) Möglichkeiten(Plural Nom) um jeden(es ein pronom hier ist und laut umwandeln, es umwandelt sich in Akkusativ Objekt) individuell nach einen Jornalist(m)  umzuwandeln(to convert. umwandeln etwas in etwas).
        
        Aus die obigen Argumente, Sozia...
        
        \bibliographystyle{plainnat}
        \bibliography{/home/donkarlo/Dropbox/repo/nd_shared_research/bibliography/bibliography.bib}
\end{document}